\usepackage[T1]{fontenc}
\usepackage[utf8]{inputenc}
\IfFileExists{french.ldf}{\usepackage[french]{babel}}{%
  \usepackage{babel}\babelprovide[import,main]{french}%
  \hyphenpenalty=10000\exhyphenpenalty=10000
}
\usepackage{microtype}
\usepackage{lmodern,amsmath,amssymb}
\usepackage[a4paper,margin=24mm,headheight=14pt]{geometry}
\usepackage{xcolor,tikz,fancyhdr,graphicx,float}
\usetikzlibrary{arrows.meta}
\usepackage{hyperref}
% Palette maîtresse. Les scripts Python lisent directement ces codes HTML.
\definecolor{accent}{HTML}{000000}
\definecolor{schemaOrange}{HTML}{FF0400}
\definecolor{schemaGreen}{HTML}{3FE61D}
\definecolor{schemaInk}{HTML}{000000}
\definecolor{schemaMuted}{HTML}{52595F}
\definecolor{schemaPurple}{HTML}{b350ff}
\definecolor{schemaNeutral}{HTML}{555555}
\definecolor{schemaGrid}{HTML}{B0B0B0}

\tikzset{
 schema/.style={font=\small,text=schemaInk,draw=schemaInk,line width=.7pt,
   >=Stealth,line cap=round,line join=round},
 schema axis/.style={draw=schemaMuted,line width=.55pt,->},
 schema curve/.style={draw=accent,line width=.95pt},
 schema vector/.style={draw=accent,line width=.9pt,->},
 schema guide/.style={draw=schemaMuted,line width=.55pt,densely dashed},
 schema approximation/.style={draw=schemaOrange,line width=.9pt},
 schema secondary/.style={draw=schemaGreen,line width=.9pt},
 schema area/.style={fill=accent!10,draw=none}
}
\hypersetup{colorlinks=true,linkcolor=accent,urlcolor=accent,pdftitle={Simulateur de forces centrales - Rappels mathématiques}}
\pagestyle{fancy}
\fancyhead{}
\fancyfoot{}
\fancyfoot[LO]{\parbox[t]{.46\textwidth}{\scriptsize\sffamily\raggedright\nouppercase{\leftmark}}}
\fancyfoot[RO]{\small\sffamily\thepage}
\fancyfoot[LE]{\small\sffamily\thepage}
\fancyfoot[RE]{\parbox[t]{.46\textwidth}{\scriptsize\sffamily\raggedleft\nouppercase{\leftmark}}}
\renewcommand{\sectionmark}[1]{\markboth{#1}{}}
\renewcommand{\subsectionmark}[1]{\markright{#1}}
\renewcommand{\headrulewidth}{0pt}
\renewcommand{\footrulewidth}{0pt}
\fancypagestyle{plain}{%
  \fancyhf{}%
  \fancyfoot[LO]{\parbox[t]{.46\textwidth}{\scriptsize\sffamily\raggedright\nouppercase{\leftmark}}}%
  \fancyfoot[RO]{\small\sffamily\thepage}%
  \fancyfoot[LE]{\small\sffamily\thepage}%
  \fancyfoot[RE]{\parbox[t]{.46\textwidth}{\scriptsize\sffamily\raggedleft\nouppercase{\leftmark}}}%
  \renewcommand{\headrulewidth}{0pt}%
  \renewcommand{\footrulewidth}{0pt}%
}
\setlength{\parindent}{0pt}
\setlength{\parskip}{4pt}
\setlength{\emergencystretch}{2em}
\setcounter{tocdepth}{1}
\numberwithin{equation}{section}
\newcommand{\sourcepage}[1]{{\small\color{gray}Source : page manuscrite #1.}\par\medskip}
\newcommand{\afaire}[1]{\par\smallskip\noindent\fbox{\parbox{0.94\linewidth}{\textbf{À compléter (manuscrit).} #1}}\par\smallskip}
\newcommand{\Em}{E_{\mathrm m}}
\newcommand{\dd}{\,\mathrm d}
\newcommand{\Ocal}{\mathcal O}
